\medskip

La population d'une espèce en voie de disparition est surveillée de près dans une réserve naturelle.

Les conditions climatiques ainsi que le braconnage font que cette population diminue de 10\,\% chaque année.

Afin de compenser ces pertes, on réintroduit dans la réserve $100$ individus à la fin de chaque année.

On souhaite étudier l'évolution de l'effectif de cette population au cours du temps. Pour cela, on modélise l'effectif de la population de l'espèce par la suite $\left(u_n\right)$ où $u_n$ représente l'effectif de la population au début de l'année $2020 + n$. 

On admet que pour tout entier naturel $n$,\: $u_n \geqslant 0$.

Au début de l'année 2020, la population étudiée compte \np{2000} individus, ainsi $u_0 = \np{2000}$.

\medskip

\begin{enumerate}
\item Justifier que la suite $\left(u_n\right)$ vérifie la relation de récurrence:

\[u_{n+1} = 0,9u_n + 100.\]

\item Calculer $u_1$ puis $u_2$.
\item Démontrer par récurrence que pour tout entier naturel $n$ :\: $\np{1000} < u_{n+1}  \leqslant u_n$.
\item La suite $\left(u_n\right)$ est-elle convergente ? Justifier la réponse.
\item On considère la suite $\left(v_n\right)$ définie pour tout entier naturel $n$ par $v_n = u_n - \np{1000}$.
	\begin{enumerate}
		\item Montrer que la suite $\left(v_n\right)$ est géométrique de raison $0,9$.
		\item En déduire que, pour tout entier naturel $n$,\: $u_n = \np{1000} \left(1 + 0,9^n\right)$.
		\item Déterminer la limite de la suite $\left(u_n\right)$.
		
En donner une interprétation dans le contexte de cet exercice.
	\end{enumerate}

\item On souhaite déterminer le nombre d'années nécessaires pour que l'effectif de la population passe en dessous d'un certain seuil $S$ (avec $S > \np{1000}$).
	\begin{enumerate}
		\item Déterminer le plus petit entier $n$ tel que $u_n \leqslant \np{1020}$.

Justifier la réponse par un calcul.
	\end{enumerate}
\begin{minipage}{0.58\linewidth}
	\begin{enumerate}[resume]
		\item Dans le programme Python ci-contre, la variable $n$ désigne le nombre d'années écoulées depuis 2020, la variable $u$ désigne l'effectif de la population. 
		
Recopier et compléter ce programme afin qu'il retourne le nombre d'années nécessaires pour que l'effectif de la population passe en dessous du seuil $S$.
	\end{enumerate}

\end{minipage}\hfill
\begin{minipage}{0.38\linewidth}
\begin{tabular}{|c l|}\hline
1&\texttt{def population(S) :}\\
2&\quad \texttt{n=0}\\
3&\quad \texttt{u=2000}\\
4& \\
5&\quad \texttt{while \ldots\ldots} :\\
6&\quad \quad \texttt{u= \ldots}\\
7&\quad \quad \texttt{n = \ldots}\\
8&\quad \texttt{return \ldots}\\ \hline
\end{tabular}
\end{minipage}
\end{enumerate}

\bigskip
